\subsubsection{nsjail}\label{nsjail}

Nsjail e\textquotesingle{} una sandbox sicura che esegue un server TCP.
Ad ogni connessione viene fatto il fork e dell\textquotesingle ambiente
e rilasciato all\textquotesingle utente per l\textquotesingle utilizzo.
L\textquotesingle ambiente separato rende impossibile agli utenti di
interferire tra di loro sebbene operino nello stesso sistema.

Nsjail e\textquotesingle{} un tool che deve essere configurato per ogni
challenge, specificando quali file dovrebbero essere esposti nel nuovo
ambiente, e quali limiti applicare.

Nsjail riesce a isolare i processi utilizzando namespaces, resource
limits ed filtri seccomp.

\subsubsection{namespaces}\label{namespaces}

\subsubsection{seccomp bpf}\label{seccomp-bpf}

Permette ad un processo di specificare un filtro da applicare alle
chiamate di sistema in uscita, specificando un comportamento (es:
uccidere il processo) quando chiamate filtrate vengono richieste dal
processo.

\begin{lstlisting}[language=C]
scmp_filter_ctx ctx;
seccomp_rule_add(ctx, ...);
seccomp_load(ctx, ...);
\end{lstlisting}

\subsubsection{cgroups}\label{cgroups}

e\textquotesingle{} un meccanismo per aggregare o partizionare i task in
gruppi gerarchici. i task sono racchiusi in sottosistemi, dove un
sottostistema si occupa di schedulare le risorse sottostando alle
limitazioni.

\passthrough{\lstinline!cgroups!} vengono utilizzati per applicare
limitazioni e fare accountig su gruppi di task.
